\documentclass[11pt, a4paper]{article}

% bfw-style.tex — gemeinsame Style-Definitionen für alle Handouts/Aufgabenhefte
% Voraussetzung: Kompilation mit xelatex (wegen fontspec)

\usepackage[ngerman]{babel}
\usepackage{geometry}
\geometry{a4paper, top=2.2cm, bottom=2.2cm, left=2.3cm, right=2.3cm}

\usepackage{fontspec}
% Fallback-Kette: Source Sans Pro -> Calibri -> Segoe UI -> Latin Modern Sans
% (Latin Modern ist immer mit MiKTeX vorhanden)
\IfFontExistsTF{Source Sans Pro}{%
  \setmainfont{Source Sans Pro}[Ligatures=TeX]%
}{\IfFontExistsTF{Calibri}{%
  \setmainfont{Calibri}[Ligatures=TeX]%
}{\IfFontExistsTF{Segoe UI}{%
  \setmainfont{Segoe UI}[Ligatures=TeX]%
}{%
  \setmainfont{Latin Modern Sans}[Ligatures=TeX]%
}}}
\IfFontExistsTF{Source Code Pro}{%
  \setmonofont{Source Code Pro}[Scale=0.92]%
}{\IfFontExistsTF{Consolas}{%
  \setmonofont{Consolas}[Scale=0.92]%
}{%
  \setmonofont{Latin Modern Mono}[Scale=0.92]%
}}

\usepackage{xcolor}
% Farbpalette: hell, freundlich, modern
\definecolor{bfwPrimary}{HTML}{0EA5A4}    % Petrol/Türkis - Primär
\definecolor{bfwPrimaryDark}{HTML}{0F766E} % dunkler Petrol für Headings
\definecolor{bfwAccent}{HTML}{F59E0B}     % Amber - Akzent
\definecolor{bfwSuccess}{HTML}{10B981}    % Grün - Erfolg / Lösung
\definecolor{bfwWarn}{HTML}{F97316}       % Orange - Achtung
\definecolor{bfwInk}{HTML}{0F172A}        % fast Schwarz
\definecolor{bfwInkSoft}{HTML}{475569}    % gedämpftes Grau
\definecolor{bfwBgSoft}{HTML}{F8FAFC}     % off-white
\definecolor{bfwBgTeal}{HTML}{ECFEFF}     % helles Türkis
\definecolor{bfwBgAmber}{HTML}{FEF3C7}    % helles Gelb
\definecolor{bfwBgRose}{HTML}{FFE4E6}     % helles Rosé für Eselsbrücken

\usepackage[colorlinks=true, linkcolor=bfwPrimaryDark, urlcolor=bfwPrimaryDark]{hyperref}
\usepackage{enumitem}
\setlist[itemize]{leftmargin=*, itemsep=2pt, topsep=2pt}
\setlist[enumerate]{leftmargin=*, itemsep=2pt, topsep=2pt}

\usepackage[most]{tcolorbox}
\usepackage{booktabs}
\usepackage{tikz}
\usetikzlibrary{decorations.pathmorphing, positioning, shapes.geometric, arrows.meta}

% Merkkästchen — wichtige Definition / Kernpunkt
\newtcolorbox{merksatz}[1][]{%
  enhanced, breakable,
  colback=bfwBgTeal,
  colframe=bfwPrimary,
  boxrule=0.6pt,
  arc=4pt,
  left=10pt, right=10pt, top=8pt, bottom=8pt,
  fonttitle=\bfseries\color{bfwPrimaryDark},
  title={\faLightbulb\ Merksatz},
  #1
}

% Eselsbrücke — Mnemoniks
\newtcolorbox{eselsbruecke}[1][]{%
  enhanced, breakable,
  colback=bfwBgRose,
  colframe=bfwAccent,
  boxrule=0.6pt,
  arc=4pt,
  left=10pt, right=10pt, top=8pt, bottom=8pt,
  fonttitle=\bfseries\color{bfwWarn},
  title={Eselsbrücke},
  #1
}

% Beispiel-Box
\newtcolorbox{beispiel}[1][]{%
  enhanced, breakable,
  colback=bfwBgSoft,
  colframe=bfwInkSoft,
  boxrule=0.4pt,
  arc=3pt,
  left=10pt, right=10pt, top=8pt, bottom=8pt,
  fonttitle=\bfseries\color{bfwInk},
  title={Beispiel},
  #1
}

% Lösungs-Box (für Aufgabenhefte)
\newtcolorbox{loesung}[1][]{%
  enhanced, breakable,
  colback=bfwBgAmber!40,
  colframe=bfwSuccess,
  boxrule=0.6pt,
  arc=3pt,
  left=10pt, right=10pt, top=8pt, bottom=8pt,
  fonttitle=\bfseries\color{bfwSuccess!70!black},
  title={Lösung},
  #1
}

% Glossar-Eintrag
\newcommand{\glossareintrag}[2]{%
  \par\noindent\textbf{\color{bfwPrimaryDark}#1} \enspace #2\par\smallskip
}

% Section-Stil — etwas farbiger als Standard
\usepackage{titlesec}
\titleformat{\section}
  {\Large\bfseries\color{bfwPrimaryDark}}
  {\thesection}{0.6em}{}
\titleformat{\subsection}
  {\large\bfseries\color{bfwPrimary}}
  {\thesubsection}{0.5em}{}
\titleformat{\subsubsection}
  {\normalsize\bfseries\color{bfwInk}}
  {\thesubsubsection}{0.4em}{}

% TikZ-Stil "skizzenhaft" — etwas weicher als Rough.js, aber stabil
\tikzset{
  roughbox/.style = {
    draw=bfwPrimaryDark, thick, fill=bfwBgTeal,
    rounded corners=4pt, inner sep=6pt
  },
  roughaccent/.style = {
    draw=bfwAccent, thick, fill=bfwBgAmber,
    rounded corners=4pt, inner sep=6pt
  },
  roughline/.style = {
    draw=bfwInkSoft, thick, -{Stealth[length=2.5mm]}
  }
}

% Bullet-Symbole (bewusst kein FontAwesome — vermeidet Font-Installations-Probleme)
\newcommand{\faLightbulb}{$\bullet$}
\newcommand{\faBookmark}{$\bullet$}

% Header/Footer
\usepackage{fancyhdr}
\pagestyle{fancy}
\fancyhf{}
\renewcommand{\headrulewidth}{0pt}
\fancyfoot[L]{\small\color{bfwInkSoft}\thepage}
\fancyfoot[R]{\small\color{bfwInkSoft} BFW M-064 Beschaffung im Büromanagement}


\title{\vspace*{-1.5cm}\Huge\bfseries\color{bfwPrimaryDark}Tag~2: Bezugsquellen, Anfragen, Angebotsvergleich\\[0.4em]
  \large\color{bfwInkSoft}\textnormal{Vom Bedarf zum besten Lieferanten --- quantitativ und qualitativ}}
\author{\textsf{Beschaffung im Büromanagement}\\
\textsf{\small\copyright\ Can Siebert\,\textperiodcentered\,2026}}
\date{12.05.2026}

\begin{document}
\maketitle
\thispagestyle{empty}

\vspace{1em}
\begin{merksatz}[title={\bfseries Worum geht es heute?}]
Nachdem Sie gestern den Beschaffungsprozess als Ganzes kennengelernt haben, gehen wir
heute die ersten vier Prozessschritte im Detail durch: \emph{Bedarfsermittlung,
Bezugsquellen­ermittlung, Anfrage und Angebotsvergleich}. Diese Schritte entscheiden,
\emph{welcher Lieferant} am Ende beauftragt wird --- und damit auch, ob Ihr Betrieb
wirtschaftlich, qualitativ und nachhaltig einkauft.
\end{merksatz}

\tableofcontents
\vspace{1em}
\hrule
\vspace{1em}

\section{Bedarfsermittlung}

\subsection{Warum die Bedarfsermittlung der wichtigste erste Schritt ist}

Jeder Beschaffungsvorgang beginnt mit der Frage: \emph{Was wird wann in welcher Menge
gebraucht?} So banal die Frage klingt, so folgenreich ist sie. Wer den Bedarf zu hoch
ansetzt, verursacht Lager- und Kapital­bindungs­kosten. Wer ihn zu niedrig ansetzt,
riskiert Versorgungs­lücken, Eil­bestellungen mit Aufpreis oder gar Produktions­stillstand.

In gut organisierten Betrieben wird die Bedarfsermittlung deshalb nicht spontan oder
,,nach Bauchgefühl'' durchgeführt, sondern auf Basis von Verbrauchsstatistiken, Stücklisten
oder Kundenaufträgen --- je nachdem, um welche Art von Material es sich handelt.

\subsection{Bedarfsarten}

Die Betriebswirtschaft unterscheidet drei Bedarfsarten nach ihrem Verwendungs­zweck:

\textbf{Primärbedarf} sind die verkaufsfertigen Endprodukte oder Dienstleistungen.
Beispiel: Ein Möbelhersteller plant für das nächste Quartal 3\,000 Schreibtische ---
das ist sein Primärbedarf.

\textbf{Sekundärbedarf} sind die Roh-, Hilfs- und Halbfertig­materialien, die in das
Endprodukt direkt eingehen. Aus den 3\,000 Schreibtischen leitet sich ab: 6\,000
Buchenholz­platten, 12\,000 Tischbeine, 24\,000 Schrauben.

\textbf{Tertiärbedarf} sind Hilfs- und Betriebsstoffe, die nicht in das Produkt eingehen,
aber für die Produktion gebraucht werden. Schmierstoff für die Säge, Reinigungs­mittel
für die Werkstatt, Druckluft, Strom.

\begin{eselsbruecke}
\textbf{P--S--T:} \emph{Primär} = Endprodukt, \emph{Sekundär} = drinsteckend,
\emph{Tertiär} = drumherum.
\end{eselsbruecke}

\subsection{Brutto- und Nettobedarf}

Während die Bedarfsarten beschreiben \emph{was} gebraucht wird, beschreibt die
Brutto-Netto-Rechnung \emph{wie viel} \emph{neu beschafft} werden muss.

\begin{merksatz}[title={Formel}]
\textbf{Nettobedarf = Bruttobedarf $-$ Lagerbestand $-$ bereits offene Bestellungen
$+$ gewünschter Sicherheits\-bestand}
\end{merksatz}

\begin{beispiel}
Sie planen für das dritte Quartal 200 Aktenordner (Bruttobedarf). Im Lager sind noch
30~Stück, weitere 50~Stück sind bereits beim Lieferanten bestellt. Der gewünschte
Sicherheits\-bestand beträgt 20~Stück.

\medskip
$\text{Nettobedarf} = 200 - 30 - 50 + 20 = \textbf{140 Stück}$

\medskip
Sie müssen also \emph{140~neue Aktenordner} beschaffen --- nicht 200, nicht 150.
\end{beispiel}

\subsection{Verbrauchsgesteuerte oder bedarfsgesteuerte Ermittlung?}

Es gibt zwei grundlegende Verfahren, den Bedarf zu ermitteln:

Bei der \emph{verbrauchsgesteuerten} Ermittlung wird der zukünftige Bedarf aus dem
\emph{Vergangenheits­verbrauch} hochgerechnet. Wenn Sie im letzten Jahr 800~Aktenordner
verbraucht haben, planen Sie 200~pro Quartal. Vorteil: einfach. Nachteil: Veränderungen
(neue Mitarbeitende, neue Kunden, geänderte Geschäfts­prozesse) werden zu spät erkannt.

Bei der \emph{bedarfsgesteuerten} Ermittlung wird der Materialbedarf konkret aus
\emph{Stücklisten} und \emph{Kundenaufträgen} berechnet. Vorteil: präzise. Nachteil:
hoher Aufwand, in der Regel nur mit IT-Unterstützung (Warenwirtschafts­system, ERP)
machbar.

In der Praxis wird kombiniert: A-Teile (siehe ABC-Analyse) bedarfsgesteuert, C-Teile
verbrauchsgesteuert.

\section{Bezugsquellen ermitteln und Lieferanten auswählen}

\subsection{Wo findet man Lieferanten?}

Bevor Sie ein Angebot einholen können, müssen Sie wissen, \emph{wen} Sie überhaupt
fragen sollen. Bezugsquellen lassen sich in interne und externe Quellen unterteilen.

\textbf{Interne Quellen} sind die hauseigene Lieferantendatei aus früheren Bestellungen,
Erfahrungs­berichte aus anderen Abteilungen sowie archivierte Angebote. Wer das pflegt,
spart bei jeder neuen Anfrage Zeit.

\textbf{Externe Quellen} sind Branchen- und Adressbücher, Fach­zeitschriften, Messen
(z.\,B.\ Frankfurter Buchmesse für Verlagsmaterial, CeBIT für IT), Online­plattformen
(Mercateo, ,,Wer liefert was''), IHK-Datenbanken sowie Empfehlungen von Geschäfts­partnern.

\subsection{Welcher Lieferant ist der richtige?}

Der billigste Lieferant ist nicht automatisch der beste. Eine seriöse Lieferanten\-auswahl
betrachtet immer mehrere Kriterien:

\begin{itemize}[itemsep=0pt]
  \item \textbf{Preis und Konditionen} --- Listenpreis, Liefererrabatt, Liefererskonto, Zahlungsziel
  \item \textbf{Qualität} --- Spezifikationstreue, Reklamationsquote, Zertifikate (z.\,B.\ ISO~9001)
  \item \textbf{Lieferzuverlässigkeit} --- Termintreue, Vollständigkeit, kommunizierte Lieferzeiten
  \item \textbf{Service} --- Reklamationsbearbeitung, Beratung, Garantie- und Gewährleistungsbedingungen
  \item \textbf{Standort und Logistik} --- Lieferzeit, Transportkosten, Versorgungssicherheit
  \item \textbf{Wirtschaftliche Stabilität} --- Bonität, Insolvenzrisiko (Schufa, Creditreform)
  \item \textbf{Nachhaltigkeit} --- Umwelt­bilanz, Sozial­standards (besonders relevant unter LkSG)
\end{itemize}

\section{ABC-Analyse}

\subsection{Das Pareto-Prinzip in der Beschaffung}

\begin{merksatz}
Die ABC-Analyse ist ein Klassifikationsverfahren, das Artikel oder Lieferanten nach
ihrem \emph{wirtschaftlichen Wertanteil} in drei Gruppen einteilt: A (sehr wichtig),
B (wichtig), C (weniger wichtig).
\end{merksatz}

Sie folgt dem Pareto-Prinzip: Etwa 20\,\% der Artikel verursachen typischerweise rund
80\,\% des Wertes. Wenn Sie diese 20\,\% identifizieren und gezielt steuern, haben Sie
den größten Hebel mit dem geringsten Aufwand.

\subsection{So führen Sie eine ABC-Analyse durch}

\begin{enumerate}
  \item Pro Artikel den Jahresverbrauchswert berechnen (Menge $\times$ Preis).
  \item Artikel nach Wert absteigend sortieren.
  \item Pro Artikel den prozentualen Anteil am Gesamtwert berechnen.
  \item Die Prozentsätze von oben nach unten kumulieren.
  \item Klassengrenzen festlegen: A bis $\approx$\,70\,\%, B bis $\approx$\,90\,\%, C der Rest.
\end{enumerate}

\begin{beispiel}
Nehmen wir 10 Bürobedarfs-Artikel mit Jahres\-verbrauchswert insgesamt 125\,200~€:

\begin{center}
\begin{tabular}{lrrrl}
\toprule
\textbf{Artikel} & \textbf{Wert} & \textbf{Anteil} & \textbf{Kumuliert} & \textbf{Klasse}\\
\midrule
A1 & 48\,000 & 38{,}3\,\% & 38{,}3\,\% & A\\
A5 & 35\,000 & 27{,}9\,\% & 66{,}2\,\% & A\\
A3 & 16\,000 & 12{,}8\,\% & 79{,}0\,\% & B\\
A7 & 9\,500 & 7{,}6\,\% & 86{,}6\,\% & B\\
A10 & 6\,350 & 5{,}1\,\% & 91{,}7\,\% & B\\
A2, A9, A4, A6, A8 & 10\,350 & 8{,}3\,\% & 100{,}0\,\% & C\\
\bottomrule
\end{tabular}
\end{center}

Ergebnis: 2 von 10 Artikeln (20\,\%) machen 66\,\% des Wertes aus --- klassisches Pareto.
\end{beispiel}

\subsection{Konsequenzen für die Beschaffungs­strategie}

\textbf{A-Teile} verdienen die höchste Aufmerksamkeit. Sie werden mit Rahmen­verträgen
gesichert, eng überwacht (kleine Sicherheits­bestände, Just-in-Time, häufige Bestell­zyklen),
und mit den Lieferanten wird strategische Partnerschaft gepflegt. Bei jedem A-Teil lohnt
sich Verhandlungs­arbeit, weil schon kleine Konditions­verbesserungen große Beträge ergeben.

\textbf{B-Teile} werden in periodischen Zyklen (z.\,B.\ monatlich) bestellt, mit mittleren
Sicherheits­beständen. Lieferanten werden regelmäßig auditiert, aber nicht so intensiv
betreut wie bei A-Teilen.

\textbf{C-Teile} werden möglichst effizient bewirtschaftet --- mit Vorrats\-beschaffung,
größeren Sicherheits­beständen und Standard­prozessen. Aufwendige Verhandlungen lohnen
sich hier nicht: Selbst wenn Sie 10\,\% beim Bleistift sparen, ist der Effekt am
Gesamt­einkauf vernachlässigbar.

\section{Anfrage}

\subsection{Begriff und rechtliche Bedeutung}

\begin{merksatz}[title={Wichtig}]
Eine \emph{Anfrage} ist rechtlich \textbf{nicht bindend}. Sie ist eine unverbindliche
Aufforderung an einen Lieferanten, ein Angebot abzugeben (lateinisch:
\emph{invitatio ad offerendum}). Erst dessen \emph{Angebot} ist ein bindender
Vertragsantrag im Sinne von \S~145 BGB.
\end{merksatz}

Das ist kein juristisches Detail, sondern hat praktische Konsequenzen: Sie können
Anfragen an mehrere Lieferanten gleichzeitig stellen, Sie können Angebote ohne Begründung
ablehnen, und Sie sind durch Ihre Anfrage zu nichts verpflichtet.

Trotzdem sollten Sie die Anfrage sorgfältig formulieren --- sie ist die Basis für das
Angebot, das Sie zurück­bekommen. Eine schlampige Anfrage führt zu schlampigen Angeboten.

\subsection{Inhaltliche Bestandteile einer Anfrage}

Eine vollständige Anfrage enthält:

\begin{itemize}[itemsep=0pt]
  \item \textbf{Genaue Artikelbezeichnung} und Spezifikation (Maße, Material, Norm, Farbe)
  \item \textbf{Menge}
  \item \textbf{Gewünschter Liefertermin}
  \item \textbf{Lieferort} (Anlieferadresse)
  \item Frage nach \textbf{Zahlungs- und Lieferbedingungen} (Skonto, Zahlungsziel, Frachtkosten)
  \item Frage nach \textbf{Verpackung}
  \item Frage nach \textbf{Garantie- bzw.\ Gewährleistungs}­konditionen
  \item \textbf{Angebotsfrist} (bis wann das Angebot vorliegen soll)
\end{itemize}

Form: Die Anfrage erfolgt schriftlich nach DIN~5008 (Geschäfts\-brief). Eine saubere
Aufmachung erhöht die Antwort­bereitschaft des Lieferanten.

\section{Angebotsvergleich}

\subsection{Quantitativ und qualitativ --- beides nötig}

Wenn drei Angebote vorliegen, muss verglichen werden. Dabei gibt es zwei
Vergleichs­ebenen, die beide notwendig sind:

\begin{itemize}[itemsep=0pt]
  \item \textbf{Quantitativer Vergleich} --- Bezugskalkulation: Welcher Lieferant ist beim Endpreis (Bezugspreis) am günstigsten?
  \item \textbf{Qualitativer Vergleich} --- Nutzwertanalyse: Welcher Lieferant erfüllt insgesamt die nicht-monetären Kriterien am besten?
\end{itemize}

Die endgültige Entscheidung kombiniert beide Ergebnisse.

\subsection{Bezugskalkulation --- das Schema}

Der Listenpreis ist nicht der Endpreis. Erst nach Berücksichtigung von Rabatten,
Skonto und Bezugskosten erhalten Sie den \emph{tatsächlichen} Preis pro Stück, den
Sie fairerweise mit anderen Angeboten vergleichen können.

\begin{merksatz}[title={Berechnungs­schema}]
\setlength{\tabcolsep}{4pt}
\begin{tabular}{ll}
& Listeneinkaufspreis\\
$-$ & Liefererrabatt\\
$=$ & Zieleinkaufspreis\\
$-$ & Liefererskonto\\
$=$ & Bareinkaufspreis\\
$+$ & Bezugskosten (Fracht, Versicherung, Verpackung)\\
$=$ & \textbf{Bezugspreis (Einstandspreis)}\\
\end{tabular}
\end{merksatz}

\begin{eselsbruecke}
\textbf{,,L $-$ R = Z, Z $-$ S = B, B + K = E''} --- Liste, Rabatt, Ziel, Skonto, Bar, Kosten, Einstand.
\end{eselsbruecke}

\begin{beispiel}
Sie bekommen drei Angebote für 100~Aktenordner:

\begin{center}
\begin{tabular}{lrrr}
\toprule
\textbf{Position} & \textbf{Lief.\,A} & \textbf{Lief.\,B} & \textbf{Lief.\,C}\\
\midrule
Listenpreis pro Stück & 4{,}50\,€ & 4{,}80\,€ & 4{,}20\,€\\
$-$ Liefererrabatt & 10\,\% & 12\,\% & 5\,\%\\
$=$ Zieleinkaufspreis & 4{,}05\,€ & 4{,}22\,€ & 3{,}99\,€\\
$-$ Liefererskonto & 2\,\% & 3\,\% & 2\,\%\\
$=$ Bareinkaufspreis & 3{,}97\,€ & 4{,}09\,€ & 3{,}91\,€\\
$+$ Bezugskosten / Stück & 0{,}25\,€ & 0{,}18\,€ & 0{,}30\,€\\
\midrule
\textbf{Bezugspreis / Stück} & \textbf{4{,}22\,€} & \textbf{4{,}27\,€} & \textbf{4{,}21\,€}\\
\bottomrule
\end{tabular}
\end{center}

Quantitativ ist Lieferant \textbf{C} mit 4{,}21\,€ am günstigsten --- knapp vor A.
Lieferant B ist trotz höchstem Rabatt am teuersten, weil sein Listenpreis am höchsten war.
\end{beispiel}

\subsection{Nutzwertanalyse --- das Schema}

Die Nutzwertanalyse erweitert den Vergleich um nicht-monetäre Kriterien.

\begin{enumerate}
  \item Kriterien festlegen (z.\,B.\ Lieferzuverlässigkeit, Qualität, Service, Nachhaltigkeit)
  \item Pro Kriterium eine Gewichtung in Prozent vergeben (Summe = 100\,\%)
  \item Pro Lieferant pro Kriterium Punkte vergeben (z.\,B.\ 1--5)
  \item Pro Kriterium: Punkte $\times$ Gewichtung = gewichteter Teilnutzen
  \item Teilnutzen aufsummieren $\to$ \emph{Nutzwert} pro Lieferant
\end{enumerate}

\begin{beispiel}
Bewertung der drei Lieferanten aus dem vorigen Beispiel:

\begin{center}
\begin{tabular}{lrrrr}
\toprule
\textbf{Kriterium} & \textbf{Gew.} & \textbf{A} & \textbf{B} & \textbf{C}\\
\midrule
Lieferzuverlässigkeit & 30\,\% & 4 $\to$ 1{,}20 & 5 $\to$ 1{,}50 & 3 $\to$ 0{,}90\\
Qualität & 25\,\% & 4 $\to$ 1{,}00 & 4 $\to$ 1{,}00 & 5 $\to$ 1{,}25\\
Service / Reklamation & 20\,\% & 3 $\to$ 0{,}60 & 4 $\to$ 0{,}80 & 4 $\to$ 0{,}80\\
Zahlungsziele & 15\,\% & 4 $\to$ 0{,}60 & 3 $\to$ 0{,}45 & 5 $\to$ 0{,}75\\
Nachhaltigkeit & 10\,\% & 3 $\to$ 0{,}30 & 5 $\to$ 0{,}50 & 2 $\to$ 0{,}20\\
\midrule
\textbf{Nutzwert gesamt} & \textbf{100\,\%} & \textbf{3{,}70} & \textbf{4{,}25} & \textbf{3{,}90}\\
\bottomrule
\end{tabular}
\end{center}

Qualitativ gewinnt Lieferant \textbf{B} klar (4{,}25 vs.\ 3{,}90 vs.\ 3{,}70).
\end{beispiel}

\subsection{Endgültige Entscheidung}

Quantitativ liegt Lieferant C vorne (4{,}21\,€), qualitativ Lieferant B (Nutzwert 4{,}25).
Bei nur 0{,}06\,€ Aufpreis gegenüber C bietet B aber deutlich höhere Lieferzuverlässigkeit
und bessere Nachhaltigkeit. Bei 100~Stück sind das 6,00\,€ Mehrkosten --- ein guter
Tausch gegen weniger Reklamations­risiko.

\begin{merksatz}
Die endgültige Lieferantenentscheidung sollte beide Vergleiche berücksichtigen ---
Bezugspreis allein führt oft zur falschen Wahl. Dokumentieren Sie Ihre Entscheidung
mit kurzer Begründung, das schützt im Streitfall.
\end{merksatz}

\section{Zusammenfassung}

\begin{enumerate}
  \item \textbf{Bedarfsarten}: Primär (Endprodukt), Sekundär (drinsteckend), Tertiär (drumherum). Brutto $-$ Lager $-$ Bestellt $+$ Sicherheit = Netto.
  \item \textbf{ABC-Analyse}: 20\,\% der Artikel = 80\,\% des Werts. A-Teile intensiv, C-Teile effizient bewirtschaften.
  \item \textbf{Anfrage}: rechtlich nicht bindend, aber inhaltlich präzise (DIN 5008).
  \item \textbf{Bezugskalkulation}: Liste $-$ Rabatt $-$ Skonto $+$ Bezugskosten = Bezugspreis.
  \item \textbf{Nutzwertanalyse}: gewichtete Bewertung qualitativer Kriterien.
  \item \textbf{Beste Wahl} = Synthese aus quantitativem und qualitativem Vergleich.
\end{enumerate}

\newpage

\section*{Glossar}
\addcontentsline{toc}{section}{Glossar}

\glossareintrag{Primärbedarf}{Verkaufsfähige Endprodukte und Dienstleistungen, die ein Unternehmen am Markt anbietet.}

\glossareintrag{Sekundärbedarf}{Roh-, Hilfs- und Halbfertig­materialien, die zur Herstellung des Primär­bedarfs benötigt werden und in das Endprodukt direkt eingehen.}

\glossareintrag{Tertiärbedarf}{Hilfs- und Betriebsstoffe, die nicht in das Endprodukt eingehen, aber für den Produktions­prozess gebraucht werden (z.\,B.\ Schmierstoffe, Reinigungsmittel).}

\glossareintrag{Bruttobedarf}{Gesamter Materialbedarf einer Periode ohne Berücksichtigung von Lagerbestand oder offenen Bestellungen.}

\glossareintrag{Nettobedarf}{Bruttobedarf abzüglich Lagerbestand und offener Bestellungen, zuzüglich Sicherheitsbestand. Das ist die Menge, die tatsächlich neu beschafft werden muss.}

\glossareintrag{Sicherheitsbestand}{Reservebestand im Lager, der Lieferengpässe oder Bedarfsschwankungen abfedert.}

\glossareintrag{Verbrauchsgesteuerte Bedarfsermittlung}{Bedarfsplanung auf Basis des Vergangen­heits­verbrauchs. Einfach, aber ungenau bei Veränderungen.}

\glossareintrag{Bedarfsgesteuerte Bedarfsermittlung}{Bedarfsplanung auf Basis von Stücklisten und konkreten Aufträgen. Präzise, aber aufwendig.}

\glossareintrag{ABC-Analyse}{Klassifikations­verfahren, das Objekte (Artikel, Lieferanten) nach ihrem wirtschaftlichen Wertanteil in drei Gruppen A (wichtig), B (mittel) und C (weniger wichtig) einteilt.}

\glossareintrag{Pareto-Prinzip}{20\,\% der Ursachen verursachen 80\,\% der Wirkung. In der Beschaffung: 20\,\% der Artikel = 80\,\% des Wertes.}

\glossareintrag{Anfrage}{Rechtlich unverbindliche Aufforderung an einen Lieferanten, ein Angebot abzugeben (\emph{invitatio ad offerendum}).}

\glossareintrag{Angebot}{Bindender Vertragsantrag eines Lieferanten nach \S\,145 BGB. Die Annahme durch den Käufer schließt den Vertrag.}

\glossareintrag{DIN 5008}{Deutsche Industrienorm für die Gestaltung von Geschäfts­briefen (Format, Anschriftenfeld, Bezugszeichenzeile).}

\glossareintrag{Bezugskalkulation}{Berechnungsschema zur Ermittlung des Bezugspreises (Einstandspreises) eines beschafften Gutes.}

\glossareintrag{Listeneinkaufspreis}{Vom Lieferanten in seiner Preisliste angegebener Bruttopreis (vor Rabatten und Skonto).}

\glossareintrag{Liefererrabatt}{Prozentualer Preisnachlass, den der Lieferant beim Listenpreis gewährt --- z.\,B.\ Mengenrabatt, Treuerabatt.}

\glossareintrag{Liefererskonto}{Prozentualer Preisnachlass bei Zahlung innerhalb einer kurzen Frist (typisch: 2--3\,\% innerhalb 10~Tagen).}

\glossareintrag{Bezugskosten}{Zusätzliche Kosten, die beim Bezug entstehen: Fracht, Versicherung, Verpackung, Zoll.}

\glossareintrag{Bezugspreis (Einstandspreis)}{Endpreis nach Abzug aller Nachlässe und Hinzurechnung der Bezugskosten. Vergleichs­basis bei der Lieferanten\-auswahl.}

\glossareintrag{Nutzwertanalyse}{Methode zur qualitativen Bewertung von Alternativen anhand gewichteter Kriterien. Auch \emph{Scoring-Modell}.}

\glossareintrag{Lieferzuverlässigkeit}{Maß für die Termin- und Mengentreue eines Lieferanten in vergangenen Lieferungen.}

\end{document}
